% 数学ノート集・共通テンプレート v2（積分論2026年度第1部の版面を参考）
% A5判。章内で定義・定理・例・注意等の番号を共通化。
% コンパイル： platex -interaction=nonstopmode math_note_template_v2.tex (2回)
%              dvipdfmx math_note_template_v2.dvi
\documentclass[a5paper,10pt,onecolumn]{jsarticle}
\usepackage[dvipdfmx,top=22mm,bottom=20mm,inner=19mm,outer=18mm,headheight=13pt,headsep=6mm]{geometry}
\usepackage{amsmath,amssymb,amsthm,mathtools}
\usepackage{enumitem,fancyhdr,titlesec}
\usepackage[dvipdfmx,bookmarks=true,colorlinks=false]{hyperref}
\usepackage{pxjahyper}

% ----- 差し替える文書情報 -----
\newcommand{\NoteTitle}{一階述語論理の健全性と完全性}
\newcommand{\NoteSeries}{数学ノート集}

% ----- 文字・組版 -----
% 参考PDFの質素な明朝＋Computer Modernの印象を引き継ぎながら，
% A5判で窮屈にならないよう本文の行間を少し確保する．
\linespread{1.13}
\setlength{\parindent}{1zw}
\setlength{\parskip}{0pt}
\setlength{\abovedisplayskip}{6pt plus 2pt minus 2pt}
\setlength{\belowdisplayskip}{6pt plus 2pt minus 2pt}
\setlength{\abovedisplayshortskip}{3pt plus 2pt minus 1pt}
\setlength{\belowdisplayshortskip}{4pt plus 2pt minus 1pt}
\setlist[enumerate]{topsep=3pt,itemsep=1pt,parsep=0pt,leftmargin=2.4zw}
\titleformat{\section}{\normalfont\normalsize\bfseries}{\thesection}{0.8zw}{}
\titlespacing*{\section}{0pt}{2.0ex plus .4ex}{.9ex}
\titleformat{\subsection}{\normalfont\normalsize\bfseries}{\thesubsection}{0.7zw}{}
\titlespacing*{\subsection}{0pt}{1.5ex plus .3ex}{.6ex}
\raggedbottom

% ----- 柱とページ番号：ページ上部、罫線なし -----
\fancypagestyle{mathnote}{%
  \fancyhf{}%
  \fancyhead[L]{\footnotesize \NoteSeries\ （\NoteTitle）}%
  \fancyhead[R]{\footnotesize \thepage}%
  \renewcommand{\headrulewidth}{0pt}%
  \renewcommand{\footrulewidth}{0pt}%
}
\pagestyle{mathnote}

% ----- 節内通し番号の定義・定理群 -----
\newtheoremstyle{mathnotetheorem}
 {4.5pt}{4.5pt}{\normalfont}{}%
 {\bfseries}{}{0.65zw}%
 {\thmname{#1}\thmnumber{ #2}\thmnote{（#3）}}
\theoremstyle{mathnotetheorem}
\newtheorem{theorem}{定理}[section]
\newtheorem{definition}[theorem]{定義}
\newtheorem{lemma}[theorem]{補題}
\newtheorem{proposition}[theorem]{命題}
\newtheorem{corollary}[theorem]{系}
\newtheorem{example}[theorem]{例}
\newtheorem{remark}[theorem]{注意}
\renewcommand{\proofname}{証明}
\renewcommand{\qedsymbol}{$\square$}
\makeatletter
\renewenvironment{proof}[1][\proofname]{%
  \par\pushQED{\qed}\normalfont\topsep4\p@\@plus2\p@\relax
  \trivlist\item[\hskip\labelsep\bfseries #1]\ignorespaces
}{\popQED\endtrivlist\@endpefalse}
\makeatother

% ----- 使うノートで適宜拡張する記号 -----
\newcommand{\R}{\mathbb{R}}
\newcommand{\norm}[1]{\lVert #1\rVert}
\newcommand{\abs}[1]{\lvert #1\rvert}

\begin{document}
\section*{読了後にできること}
本稿を読むための前提は，一階論理の構文・意味論，証明体系とモデル．
本稿の具体的な到達課題は，自然演繹の証明とモデルでの真理を区別し，具体式の恒真性を導く．

\section{言語・項・論理式}

本稿では古典的一階論理を扱う．
言語は可算で，等号を含むと仮定する．
証明体系としてHilbert型の体系を固定する．

\begin{definition}
\emph{一階言語} $\mathcal L$ は，
定数記号，それぞれ有限個の引数を持つ関数記号と
関係記号の集合からなる．
変数記号 $x_0,x_1,\dots$ を用いる．
\emph{項}は変数・定数から出発し，
関数記号の適用を有限回繰り返して得る式である．
\emph{原子論理式}は $t=s$ または
$R(t_1,\dots,t_n)$ の形の式である．
論理式は原子論理式から
$\neg,\land,\lor,\to,\forall,\exists$ によって
有限回の構成を繰り返して得る式とする．
\end{definition}

\begin{definition}
構造 $\mathcal M$ とは，空でない集合 $M$ と，
各定数記号 $c$ に対する元 $c^{\mathcal M}\in M$，
各関数記号 $f$ に対する演算 $f^{\mathcal M}$，
各関係記号 $R$ に対する部分集合 $R^{\mathcal M}$ を
指定したものをいう．
変数割当て $s:\{x_i\}\to M$ のもとで項の値
$t^{\mathcal M}[s]$ を再帰的に定める．
\end{definition}

\begin{definition}
満足関係 $\mathcal M\models\varphi[s]$ は
論理式の構成に関する帰納法で定義する．
原子論理式については実際の等号・関係を用い，
命題結合子には通常の真理値を用いる．
量化について
\[
 \mathcal M\models\exists x\,\varphi[s]
 \iff
 \text{ある }a\in M\text{ に対し }
 \mathcal M\models\varphi[s(x\mapsto a)]
\]
とし，$\forall$ は全ての $a\in M$ に対する成立とする．
\end{definition}

\begin{lemma}[代入補題]\label{lem:substitution}
項 $t$ が論理式 $\varphi$ の変数 $x$ に
自由に代入可能であるとする．
このとき
\[
 \mathcal M\models\varphi[t/x][s]
 \iff
 \mathcal M\models\varphi[s(x\mapsto t^{\mathcal M}[s])]
\]
が成り立つ．
\end{lemma}
\begin{proof}
原子論理式の場合，項の評価が関数の構成に関して
帰納的に定義されることから直接成立する．
否定と命題結合子については，
帰納法の仮定を各部分論理式へ適用すればよい．
量化 $\exists y\,\psi$ の場合，
$y\ne x$ なら $t$ は $y$ の量化に捕獲されないよう
変数名を変更してよい．
すると
\[
 (\exists y\,\psi)[t/x]=\exists y\,(\psi[t/x])
\]
である．
左辺の真理値はある $a\in M$ に対して
$\psi[t/x]$ が割当て $s(y\mapsto a)$ のもとで成立することと
同値である．
$t$ は $y$ を自由変数として含まない
（そうでなければ代入可能性に反する）ので，
$t^{\mathcal M}[s(y\mapsto a)]=t^{\mathcal M}[s]$ である．
帰納法の仮定を適用すれば，
この条件は右辺の満足関係と同値となる．
$y=x$ の場合は量化された $x$ は自由でないので
置換によって式が変わらず，主張は明らかである．
全称量化も同様である．
\end{proof}

\section{Hilbert型の証明体系}

\begin{definition}
$\Gamma$ を論理式の集合とする．
$\Gamma$ から $\varphi$ が\emph{形式的に証明可能}であるとは，
有限列 $\varphi_1,\dots,\varphi_n=\varphi$ があり，
各項が $\Gamma$ の元，以下の公理，
または先行項から推論規則で得られることをいう．
\begin{enumerate}
\item 全ての命題論理の恒真式の代入例．
\item $\forall x\,\psi\to\psi[t/x]$．
\item $\forall x(\alpha\to\beta)\to
 (\alpha\to\forall x\,\beta)$
（$x$ が $\alpha$ に自由に現れないとき）．
\item 等号の反射律 $t=t$ と，
等しい項を原子式の引数内で置換するための
合同公理．
\end{enumerate}
推論規則はmodus ponensと全称化である．
全称化は，その導出で未解消の仮定に自由に現れない変数
についてのみ行う．
$\Gamma\vdash\varphi$ と書く．
\end{definition}

\begin{remark}
等号の合同公理は，
全ての関数記号 $f$ について
\[
 t_i=s_i\ (1\le i\le n)\ \Longrightarrow\
 f(t_1,\dots,t_n)=f(s_1,\dots,s_n)
\]
および関係記号 $R$ について
\[
 t_i=s_i\ (1\le i\le n),\ R(t_1,\dots,t_n)
 \ \Longrightarrow\ R(s_1,\dots,s_n)
\]
を形式化した公理図式である．
\end{remark}

\begin{theorem}[健全性定理]\label{thm:soundness}
$\Gamma\vdash\varphi$ なら
$\Gamma\models\varphi$ が成り立つ．
\end{theorem}
\begin{proof}
$\Gamma$ を満たす任意の構造 $\mathcal M$ と割当て $s$ を取る．
証明列に関する帰納法を行う．
仮定として採用した論理式は
$\mathcal M$ で真である．
命題論理公理は真理値表によって常に真である．
全称量化の例化公理は代入補題\ref{lem:substitution}から真であり，
量化子の分配公理は
$x$ が $\alpha$ に自由に現れないという条件から真である．
等号の公理は等号の解釈が実際の等号であり，
関数・関係が同じ引数値に同じ結果を与えることから真である．
modus ponensは真な $\alpha$ と
$\alpha\to\beta$ から真な $\beta$ を導く．
全称化については，量化する変数は未解消の仮定に
自由に現れないから，その変数の割当てを任意に変更しても
全ての前提は真である．
従って帰納法の仮定によって，その変数の全ての値で
結論は真となり，全称化した式も真となる．
よって証明列の全ての式は前提を満たす構造で真であり，
最後の式 $\varphi$ も真となる．
\end{proof}

\begin{lemma}[演繹定理]\label{lem:deduction}
$\Gamma\cup\{\alpha\}\vdash\beta$ なら
$\Gamma\vdash\alpha\to\beta$ が成り立つ．
\end{lemma}
\begin{proof}
$\Gamma\cup\{\alpha\}$ からの証明列に関する帰納法で，
各 $\varphi_i$ に対して
$\Gamma\vdash\alpha\to\varphi_i$ を示す．
$\varphi_i=\alpha$ の場合は命題論理の恒真式
$\alpha\to\alpha$ を用いる．
公理または $\Gamma$ の元なら，
$\varphi_i\to(\alpha\to\varphi_i)$ が恒真式なので従う．
modus ponensの場合，
$\alpha\to\psi$ と
$\alpha\to(\psi\to\varphi_i)$ から
命題論理の恒真式により $\alpha\to\varphi_i$ を導ける．
全称化の場合は，変数 $x$ は未解消の仮定 $\alpha$ に
自由に現れない．
帰納法の仮定 $\Gamma\vdash\alpha\to\psi$ を
全称化して $\Gamma\vdash\forall x(\alpha\to\psi)$ とし，
量化子の分配公理から
$\Gamma\vdash\alpha\to\forall x\,\psi$ を得る．
最後の式に適用すれば結論が従う．
\end{proof}

\section{無矛盾性と最大無矛盾理論}

\begin{definition}
論理式集合 $\Gamma$ が\emph{無矛盾}であるとは，
ある論理式 $\chi$ について
$\Gamma\vdash\chi$ と
$\Gamma\vdash\neg\chi$ が同時に成立しないことをいう．
$\Gamma$ が完全であるとは，
任意の文 $\sigma$ について
$\Gamma\vdash\sigma$ または
$\Gamma\vdash\neg\sigma$ が成り立つことをいう．
\end{definition}

\begin{lemma}[無矛盾性の拡張]\label{lem:maximal}
可算言語における無矛盾な文の集合 $\Gamma$ は，
無矛盾かつ完全な文の集合 $\Gamma^*$ へ拡張できる．
\end{lemma}
\begin{proof}
言語の文を $\sigma_0,\sigma_1,\dots$ と列挙する．
$\Gamma_0=\Gamma$ と置き，
$\Gamma_n\cup\{\sigma_n\}$ が無矛盾なら
$\Gamma_{n+1}=\Gamma_n\cup\{\sigma_n\}$，
そうでなければ
$\Gamma_{n+1}=\Gamma_n\cup\{\neg\sigma_n\}$
と定める．
後者の選択でも無矛盾であることを示す．
もし両方の追加で矛盾が導けるなら，
演繹定理\ref{lem:deduction}と古典命題論理の法則から
$\Gamma_n\vdash\sigma_n$ および
$\Gamma_n\vdash\neg\sigma_n$ が従い，
$\Gamma_n$ の無矛盾性に反する．
従って帰納的に各 $\Gamma_n$ は無矛盾である．
$\Gamma^*=\bigcup_n\Gamma_n$ と置く．
証明は有限個の前提しか使用しないので，
もし $\Gamma^*$ が矛盾するなら
ある $\Gamma_N$ で既に矛盾するはずである．
従って $\Gamma^*$ は無矛盾である．
全ての文 $\sigma_n$ について，その文または否定が
$\Gamma^*$ に属するので完全である．
\end{proof}

\section{Henkin定数と証人}

\begin{lemma}[証人定数の追加]\label{lem:henkin-step}
$\Gamma$ を無矛盾な文の集合，
$\exists x\,\varphi(x)$ を文とし，
$c$ を $\Gamma$ と $\varphi$ に現れない新しい定数記号とする．
このとき
\[
 \Gamma\cup\{\exists x\,\varphi(x)\to\varphi(c)\}
\]
は無矛盾である．
\end{lemma}
\begin{proof}
矛盾すると仮定する．
演繹定理と命題論理の恒真式によって
\[
 \Gamma\vdash\exists x\,\varphi(x),
 \qquad
 \Gamma\vdash\neg\varphi(c)
\]
となる．
$c$ は元の言語に現れない新しい記号なので，
後者の導出で $c$ を新しい変数 $y$ に一様に置き換えられる．
従って $\Gamma\vdash\neg\varphi(y)$ を得る．
$y$ は元の仮定に自由に現れないから全称化が許され，
$\Gamma\vdash\forall y\,\neg\varphi(y)$ となる．
$\neg\exists x\,\varphi(x)$ は
$\forall x\,\neg\varphi(x)$ と同値であるため，
$\Gamma$ が矛盾し仮定に反する．
従って追加後も無矛盾である．
\end{proof}

\begin{theorem}[Henkin拡張]\label{thm:henkin}
可算言語の無矛盾な文集合 $\Gamma$ は，
可算個の新しい定数を加えた言語の
無矛盾かつ完全な理論 $T$ に拡張でき，
全ての存在文 $\exists x\,\varphi(x)$ に対し，
ある定数 $c_\varphi$ について
\[
 T\vdash\exists x\,\varphi(x)\to\varphi(c_\varphi)
\]
が成り立つようにできる．
\end{theorem}
\begin{proof}
まず元の言語に可算無限個の新しい定数を加え，
各存在文 $\exists x\,\varphi(x)$ を順に列挙する．
各段階で未使用の定数 $c$ を選び，
補題\ref{lem:henkin-step}による証人公理を追加する．
この操作を一度行うだけでは，新しい定数を含む存在文が
後から生じるため，言語の拡張と証人の追加を
可算回繰り返す．
各有限段階は補題\ref{lem:henkin-step}により無矛盾であり，
証明の有限性から全ての段階の和集合も無矛盾である．
さらに補題\ref{lem:maximal}を適用し，
全ての文を決定する完全な無矛盾拡張 $T$ を取る．
証人公理は拡張で保たれるため，
$T$ は全ての存在文に証人定数を持つ．
\end{proof}

\section{項モデルと真理補題}

\begin{definition}
Henkin理論 $T$ の言語において，
閉項の集合 $\mathrm{Term}_0$ 上に
\[
 t\sim s\iff T\vdash t=s
\]
と定める．
閉項の同値類全体を $M_T$ とする．
関数記号 $f$ の解釈を
\[
 f^{M_T}([t_1],\dots,[t_n])
 =[f(t_1,\dots,t_n)]
\]
と定め，関係記号 $R$ について
\[
 M_T\models R([t_1],\dots,[t_n])
 \iff T\vdash R(t_1,\dots,t_n)
\]
と定める．
これを\emph{Henkin項モデル}という．
\end{definition}

\begin{lemma}[項モデルのwell-defined性]\label{lem:term-model}
$\sim$ は同値関係であり，
項モデルの演算と関係の解釈は代表元に依存しない．
また $M_T$ は空でない．
\end{lemma}
\begin{proof}
等号の公理から反射律 $t=t$ が証明可能である．
対称律と推移律も等号の合同公理と命題論理の推論によって
導けるから $\sim$ は同値関係である．
$t_i\sim s_i$ なら等号の合同公理から
\[
 T\vdash f(t_1,\dots,t_n)=f(s_1,\dots,s_n)
\]
が成立し，同じ項同値類を与える．
関係記号の解釈についても，
等号の合同公理により
$T\vdash R(t_1,\dots,t_n)$ と
$T\vdash R(s_1,\dots,s_n)$ は同値である．
従ってwell-definedである．
新しい定数記号が少なくとも一つ存在するため，
閉項の同値類があり $M_T$ は空でない．
\end{proof}

\begin{lemma}[真理補題]\label{lem:truth}
$T$ を無矛盾・完全・Henkinな理論とする．
任意の文 $\varphi$ について
\[
 M_T\models\varphi
 \quad\Longleftrightarrow\quad T\vdash\varphi
\]
が成り立つ．
\end{lemma}
\begin{proof}
論理式の構成に関する帰納法で証明する．
閉項の解釈は項同値類そのものなので，
等号の原子式については $\sim$ の定義から成立する．
関係記号の原子式については項モデルの定義から成立する．
否定 $\neg\psi$ については，
$T$ が完全かつ無矛盾であるから
\[
 T\vdash\neg\psi\iff T\not\vdash\psi
\]
が成り立つ．
帰納法の仮定と古典論理の否定の意味論を比較すればよい．
連言については
$T\vdash\psi\land\chi$ が
$T\vdash\psi$ かつ $T\vdash\chi$ と同値であり，
意味論でも同様なので成立する．
他の命題結合子も同様である．

存在量化 $\exists x\,\psi(x)$ を考える．
$T\vdash\exists x\,\psi(x)$ なら，
Henkin性によって
$T\vdash\psi(c)$ となる定数 $c$ が存在する．
帰納法の仮定から
$M_T\models\psi([c])$ なので
$M_T\models\exists x\,\psi(x)$ となる．
逆にモデルで存在文が真なら，
ある閉項同値類 $[t]\in M_T$ が存在して
$M_T\models\psi([t])$ となる．
帰納法の仮定から $T\vdash\psi(t)$ であり，
存在量化への導入が証明可能なので
$T\vdash\exists x\,\psi(x)$ を得る．
全称量化は否定と存在量化で定義できるため，
これで全ての論理式について帰納法が完了する．
\end{proof}

\section{完全性定理}

\begin{theorem}[Gödelの完全性定理]\label{thm:completeness}
可算一階言語の文集合 $\Gamma$ と文 $\varphi$ に対し
\[
 \Gamma\models\varphi
 \quad\Longleftrightarrow\quad
 \Gamma\vdash\varphi
\]
が成り立つ．
特に無矛盾な一階理論はモデルを持つ．
\end{theorem}
\begin{proof}
右から左は健全性定理\ref{thm:soundness}による．
逆に $\Gamma\models\varphi$ かつ
$\Gamma\not\vdash\varphi$ と仮定する．
もし $\Gamma\cup\{\neg\varphi\}$ が矛盾すれば，
演繹定理\ref{lem:deduction}と古典命題論理の法則から
$\Gamma\vdash\varphi$ が従うので，
$\Gamma\cup\{\neg\varphi\}$ は無矛盾である．
定理\ref{thm:henkin}により，
この集合を含む無矛盾・完全・Henkin理論 $T$ が存在する．
補題\ref{lem:truth}からその項モデル $M_T$ は
$T$ の全ての文を満たす．
従って $M_T\models\Gamma$ かつ
$M_T\models\neg\varphi$ であり，
$\Gamma\models\varphi$ に矛盾する．
ゆえに $\Gamma\vdash\varphi$ である．

最後に無矛盾な $\Gamma$ に
Henkin拡張と項モデルを適用すれば
$\Gamma$ を満たすモデルを得られる．
\end{proof}

\begin{remark}
ここでの「完全性」とは，意味論的に全てのモデルで真となる
一階の帰結が形式的に証明可能であるという意味である．
これは，十分強い再帰的公理化可能な算術理論が
全ての算術文を決定できないという
Gödelの不完全性定理とは異なる主張である．
\end{remark}

\section{到達課題の完全解答}
\noindent\textbf{量化式の妥当性を証明する．}
\begin{proposition}
一階論理の任意の構造 $\mathcal M$ と述語 $P$ に対して
\[
(\forall x\,P(x))\longrightarrow(\exists x\,P(x))
\]
は，$\mathcal M$ の台集合が非空であれば妥当である．
\end{proposition}
\begin{proof}
$\mathcal M\models\forall xP(x)$ と仮定する．
非空性から元 $a$ を一つ選べる．
全称量化の意味から $\mathcal M\models P(a)$ であり，
存在量化の意味から $\mathcal M\models\exists xP(x)$．
したがって含意は真である．
\end{proof}
\begin{proposition}
$(\exists x\,\forall y\,R(x,y))\to
(\forall y\,\exists x\,R(x,y))$ は妥当である．
\end{proposition}
\begin{proof}
前件が真ならある元 $a$ が全ての $y$ に対して
$R(a,y)$ を満たす．任意の $y$ に対し存在量化の
証人として同じ $a$ を選べるので後件が真となる．
\end{proof}
\noindent\textbf{到達確認．}
$(\forall y\,\exists x\,R(x,y))\to
(\exists x\,\forall y\,R(x,y))$ は妥当か．
\par\noindent\textbf{解答．}
一般には妥当でない．台集合を $\{0,1\}$ とし，
$R(x,y)$ を $x=y$ と解釈する．
各 $y$ に対して $x=y$ を選べるから前件は真だが，
全ての $y$ に等しい固定した $x$ は存在しない．
\end{document}